\section{Extension to General Bilinear Games}\label{h1:linear-games}

In this section, we return to the general bilinear setting, where the action sets $\calX \subset \dR^{d_x}, \calY \subset \dR^{d_y}$ are arbitrary finite sets with large cardinalities $m_x$ and $m_y$.
We assume without loss of generality that $\calX$ and $\calY$ span $\dR^{d_x}$ and $\dR^{d_y}$, respectively. We further assume that $\normm{x}_2\leq 1$, $\normm{y}_2\leq 1$ for any $x\in \calX$ and $y\in \calY$, and $\normm{A}_2 \leq 1$; this implies $u(x,y)=\inner{x}{Ay}\in [-1, 1]$.
This framework generalizes the stochastic linear bandit problem \citep{abbasi-yadkori2011improved}, which is recovered when $\calY$ is a singleton ($|\calY|=1$).
Similar to the motivation behind linear bandits, we investigate this setting to derive
regret bounds with dominant terms that scale polynomially with dimensions $d_x, d_y$, but only logarithmically with the action set sizes $m_x,m_y$.

\subsection{Uninformed Learning in Bilinear Games with \TsallisSPM{}}\label{h2:tsallis-spm}

As Tsallis-INF is shown to be successful in normal-form games (\Cref{h2:tsallis-inf}), it is more than natural to extend it to bilinear games.
Indeed, for adversarial linear bandits, \citet{ito2024adaptive} have generalized Tsallis-INF and proposed an algorithm, which we will refer to as \TsallisSPM, to achieve similar best-of-both-worlds results.
We consider using the same algorithm for uninformed learning in general bilinear games, shown in \Cref{alg:tsallis-spm}.
It is again based on the FTRL framework,
but uses a variant of the Tsallis entropy regularizer (\Cref{line:FTRL}),
a reward estimator tailored to the linear structure (\Cref{line:estimator}), and a more sophisticated exploration and learning rate schedule (\Cref{line:explore} and \Cref{line:learning_rate}).

\begin{algorithm}[t]
\caption{\TsallisSPMLong}\label{alg:tsallis-spm}
\begin{algorithmic}[1]
\Require A Tsallis entropy parameter $\alpha\in (0, 1)$ (\textit{c.f.} \eqref{eq:tsallis}), learning rate parameters $\beta_1>0$, $\bar{\beta}\geq 0$, a fixed exploration distribution $\pzero$ with a variance ratio $c>0$.
\For{$t=1,2,\dots$}
  \State Let $\phat!t = \argmax_{p\hspace{0.08em}\in \simplex^{m_x}} \cbrm[\big]{\inner{p}{\sum_{s<t} g!s} + \beta_t \varphi_\alpha(p)+\bar{\beta} \varphi_{1-\alpha}(p)}.$
  \State Set $z_t=\frac{d_x\rbr{\min\{\normm{\phat!t}_\infty,1-\normm{\phat!t}_\infty\}}^{1-\alpha}}{1-\alpha}$.\label{line:FTRL}
  \State Set the exploration coefficient $\gamma_t=\frac{4c z_t}{\beta_t}$.
  \State Compute the distribution $p!t = (1-\gamma_t) \phat!t + \gamma_t \pzero$. \label{line:explore}
  \State Sample $x!t\in \calX$ from the distribution specified by $p!t$ and play it.
  \State Observe the reward $r_t$.
  \State Define reward estimator $g!t\in \dR^\calX$, where
      $g!t_x = r_t \inner{x!t}{S(p!t)^{-1}x}$
      and $S(p)\defeq \EE_{x\sim p}[xx^\trans]$. \label{line:estimator}
  \State Set the learning rate of the next round $\beta_{t+1}=\beta_t+\frac{z_t}{\beta_t h_t}$ where $h_t=\varphi_a(\phat!t)$. \label{line:learning_rate}
\EndFor
\end{algorithmic}
\end{algorithm}

While the algorithm is from prior work, we again provide a novel analysis for the game setting, using the same instance-dependent parameters $\DRmin$, $\DCmin$, $\DeltaMix$ defined in \Cref{h2:tsallis-inf} (note that in these definitions, $\calX$ and $\calY$ are now general sets instead of the standard basis sets.). %

\begin{theorem}\label{thm:tsallis-bilinear}
    With $\alpha=1-\frac{1}{4\log m_x}$, $\beta_1= \frac{8 c d_x}{1-\alpha}$, $\bar{\beta}= \frac{32d_x}{(1-\alpha)^2 \beta_1}$, and a suitable choice of $\pzero$, in any bilinear game, against any adaptive adversary, \TsallisSPM{} achieves a worst-case regret bound of $\psmr_T=\order\rbrm[\big]{\sqrt{T d_x \log m_x} + m_x \log m_x}$.
    Additionally,
    \begin{itemize}[noitemsep,topsep=0pt]
        \item If the game has a strict PSNE, we further have $\psmr_T=\order\rbrm[\big]{\frac{d_x \log m_x}{\DRmin\DCmin}\log T + m_x \log m_x}$.
        \item If the game has no PSNEs, we further have $\psmr_T = \order\rbrm[\big]{\frac{d_x \log m_x}{\DeltaMix} + m_x \log m_x}$.
    \end{itemize}
\end{theorem}


The proof is similar to \Cref{thm:tsallis} and is deferred to \Cref{h1a:proof-tsallis-bilinear}.
Importantly, all $T$-dependent terms have only logarithmic dependence on $m_x$.
The additive $m_x\log m_x$ term is not ideal but was present in~\citet{ito2024adaptive} already.
On the other hand, if one were to ignore the linear structure, treat this as a normal game with $m_x$ actions (for the learner), and apply Tsallis-INF, then the bound would have a polynomial dependency, as either $\sqrt{m_x T}$ (worst-case) or $\frac{1}{\DCmin}m_x\log T$ (instance-dependent).

%

\subsection{Informed Learning in Bilinear Games with \PureLinUCB{}}\label{h2:linucb}

Finally, we turn to the informed feedback setting for general bilinear games.
Analogous to our approach in \Cref{h1:pure-ucb}, we extend the classical LinUCB algorithm \citep{abbasi-yadkori2011improved} to the game-theoretic setting.





The algorithm is based on the observation that the reward of a bilinear game is linear in its matrix $A$: $u(x,y)=\inner{x}{Ay}=\inner{A}{x y^\trans}$. The feedback $r_t$ in each round is therefore a noisy linear projection of $A$. 
The algorithm can thus maintain an estimate for $A$, utilize a LinUCB-style confidence ellipsoid to determine an optimistic reward estimate for each action pair, and play the pure maximin action based on the upper confidence bounds, analogous to \PureUCB.
The formal description of this algorithm, which we name \PureLinUCB{}, is deferred to \Cref{h1a:proof-purelinucb}.


To analyze this algorithm, we let $\DeltaLin>0$ be such that $\vStar - u(x,y)\in (-\infty,0]\cup [\DeltaLin, +\infty)$. This parameter always exists since $\calX$ and $\calY$ are finite. In the special case of normal-form games, the parameter $\DeltaLin$ would be equal to $\min_{x,y} \cbr{\Delta_{xy}\mid \Delta_{xy}>0}$. 

We show that \PureLinUCB{} effectively exploits the low-dimensional nature of the bilinear game.
The full analysis is deferred to \Cref{h1a:proof-purelinucb}, with
the main result summarized below.

\begin{theorem}\label{thm:purelinucb-informal}
    For any bilinear game, against any adaptive adversary, \PureLinUCB{} achieves a worst-case regret of $\psmr_T=\order\rbrm[\big]{d_x d_y \sqrt{T} \log T}$ and an instance-dependent regret of $\psmr_T=\order\rbrm[\big]{\frac{1}{\DeltaLin}d_x^2 d_y^2 \log^2 T}$ simultaneously.
\end{theorem}

Note that the bounds scale polynomially with the feature dimensions $d_x$ and $d_y$, but have no dependence on $m_x$ and $m_y$ at all (unlike \Cref{thm:tsallis-bilinear}).
Once again, if one were to treat this as a normal-form game and apply \PureUCB{}, then its bound suffers a polynomial dependency on $m_x$ and $m_y$.
